\section{A symmetric two-binary variant}\label{app:symmetric}
In the instance of Theorem~\ref{thm:lower}, the improving residuals near
zero are all positive, so a positive multiplier halves the required
coefficient. This appendix gives a variant in which residuals of both
signs occur at cost $-1$. Any nonzero multiplier then hurts one of the two
sides, the optimal multiplier is zero, and the optimized and
zero-multiplier thresholds coincide at $2^{2^d}$.

Keep the chain $A_d$ and $\delta_d=2^{-2^d}$ of~\eqref{eq:chain}, the
objective $-q$ and the relaxed equality $y=0$. Add a second binary
variable $b\in\{0,1\}$ and replace the switching inequality
in~\eqref{eq:onebinary} by the pair
\begin{equation}\label{eq:symmetric}
 y\ge a_d-(1-q)-2(1-b),\qquad
 y\le-a_d+(1-q)+2b.
\end{equation}
The native set is thus
$\{(a,y,q,b):a\in A_d,\ -1\le y\le1,\ q,b\in\{0,1\},\ \eqref{eq:symmetric}\}$.
At $q=1$, the choice $b=1$ forces $y\ge a_d$ and the choice $b=0$ forces
$y\le-a_d$, so the improving slices lie on both sides of $y=0$. We write
$L^{\rm sym}_{d,\rho}(\lambda)$ and $D^{\rm sym}_{d,\rho}$ for the
penalized and dual values with the absolute-value penalty.

\begin{proposition}\label{prop:symmetric}
Let $d\ge1$. The variant has optimal value $0$. For every $\rho\ge0$ and
$\lambda\in\R$, with $c=\rho-|\lambda|$,
\begin{equation}\label{eq:symdual}
 L^{\rm sym}_{d,\rho}(\lambda)=\min\{0,\,-1+c\delta_d,\,-1+c\},\qquad
 D^{\rm sym}_{d,\rho}=\min\{0,\,-1+\rho\delta_d\},
\end{equation}
and $\lambda=0$ attains the dual value. Consequently:
\begin{enumerate}
\item[(a)] the least penalty and the zero-multiplier threshold both equal
$\delta_d^{-1}=2^{2^d}$, which has $2^d+1$ bits, and at
$\rho=\delta_d^{-1}$ the multiplier $\lambda=0$ is the only one at which
$\rho$ is exact;
\item[(b)] for $0\le\epsilon<1$, the gap $0-D^{\rm sym}_{d,\rho}$ is at
most $\epsilon$ exactly when $\rho\ge(1-\epsilon)\delta_d^{-1}$;
\item[(c)] every $\rho>\delta_d^{-1}$ is minimizer-exact at multiplier
zero;
\item[(d)] every integer slice $(q,b)$ contains a point at which every
native inequality in the continuous variables holds with slack at least
$1/8$; for $q=0$ this point also satisfies $y=0$.
\end{enumerate}
\end{proposition}
\begin{proof}
As for Theorem~\ref{thm:lower}, we first compute the projection onto
$(q,y)$, on which the objective and the residual depend, and then minimize
over it.

\emph{Projection.} Fix $a\in A_d$. Using $-1\le y\le1$ and
$\delta_d\le a_d\le1$, the four slices give the following $y$-intervals:
$[a_d,1]$ for $(q,b)=(1,1)$, since the upper bound $2-a_d$ is at least
one; $[-1,-a_d]$ for $(1,0)$, since the lower bound $a_d-2$ is at most
$-1$; $[a_d-1,1]$ for $(0,1)$; and $[-1,1-a_d]$ for $(0,0)$. At $q=1$
every native point therefore has $|y|\ge a_d\ge\delta_d$. At $a=\bar a$
the two intervals at $q=1$ are $[\delta_d,1]$ and $[-1,-\delta_d]$, and
the two intervals at $q=0$ are $[\delta_d-1,1]$ and $[-1,1-\delta_d]$,
which cover $[-1,1]$. Hence the projection is
\[
 (\{0\}\times[-1,1])\ \cup\
 (\{1\}\times([-1,-\delta_d]\cup[\delta_d,1])).
\]
Since $\delta_d>0$, the relaxed equality $y=0$ forces $q=0$, and the
optimal value is $0$.

\emph{Penalized value.} For a point with residual $y$, the term
$\lambda y+\rho|y|$ is at least $c|y|$, with equality for one sign of $y$.
Both signs occur on each part of the projection, and each part is
symmetric in $y$, so the minimum of $\lambda y+\rho|y|$ over a part equals
the minimum of $cs$ over the attainable magnitudes $s=|y|$. At $q=0$ these
are $s\in[0,1]$, giving $\min\{0,c\}$. At $q=1$ they are
$s\in[\delta_d,1]$, giving $-1+\min\{c\delta_d,c\}$. The term $c$ is
larger than $-1+c$, so it can be dropped, which proves the first formula
in~\eqref{eq:symdual}. This expression is nondecreasing in $c$, and
$c\le\rho$ with equality only at $\lambda=0$. Hence the dual value is
attained at $\lambda=0$ and equals $\min\{0,-1+\rho\delta_d,-1+\rho\}$.
Since $\delta_d\le1$, the last term can be dropped, which proves the
second formula.

\emph{Parts (a)--(c).} By~\eqref{eq:symdual}, $D^{\rm sym}_{d,\rho}=0$
exactly when $\rho\ge\delta_d^{-1}$, and the same holds for
$L^{\rm sym}_{d,\rho}(0)$. At $\rho=\delta_d^{-1}$ and $\lambda\ne0$ we
have $c<\delta_d^{-1}$, so $-1+c\delta_d<0$. The integer $2^{2^d}$ has
$2^d+1$ bits. For (b), the gap is $\max\{0,1-\rho\delta_d\}$, which is at
most $\epsilon\ge0$ exactly when $\rho\ge(1-\epsilon)\delta_d^{-1}$. For
(c), take $\lambda=0$ and $\rho>\delta_d^{-1}$. A native point with $q=1$
has penalized value $-1+\rho|y|\ge-1+\rho\delta_d>0$, and a native point
with $q=0$ and $y\ne0$ has penalized value $\rho|y|>0$. So the minimizers
are exactly the feasible points.

\emph{Part (d).} Set $a_i=3/8$ for all $i$. As in the proof of
Theorem~\ref{thm:lower}, the chain and its bounds have slack at least
$1/8$. On the slices with $q=0$ take $y=0$: the inequalities
in~\eqref{eq:symmetric} have slacks $21/8$ and $5/8$ for $b=0$, and $5/8$
and $21/8$ for $b=1$. On the slices with $q=1$ take $y=(2b-1)/2$: for
$b=1$ the slacks are $1/8$ and $9/8$, and for $b=0$ they are $9/8$ and
$1/8$. The bounds on $y$ have slack at least $1/2$. The smallest slack is
$1/8$.
\end{proof}

At $\rho=\delta_d^{-1}$ and $\lambda=0$, the points $(q,y)=(1,\pm\delta_d)$
have penalized value $0$, so this coefficient is exact but not
minimizer-exact. The variant has two binary variables, $d+1$ continuous
variables, a linear objective and the same chain of $d-1$ convex
quadratic inequalities as the instance of Theorem~\ref{thm:lower}. The
fractional-power augmentation of Section~\ref{sec:lower-scope} behaves
as before on this variant, because $|y|\ge\delta_d$ on its slices with
$q=1$: the coefficient two is exact, and every larger coefficient is
minimizer-exact at multiplier zero.
